\documentclass[main.tex]{subfiles}

\begin{document}
	
\subsection{General Data}	
	
Pri získavaní údajov z platformy DataCube (\url{https://datacube.statistics.sk/}) sme sa zamerali na zhromažďovanie čo najväčšieho množstva relevantných dát, ktoré by mohli slúžiť ako indikátory pre pochopenie rozhodovania slovenských voličov v parlamentných voľbách. 

Našou hlavnou myšlienkou bolo identifikovať faktory, ktoré najvýraznejšie ovplyvňujú volebnú podporu politických strán. Tieto údaje by mali slúžiť ako základ pre modelovanie, ktoré nám umožní predikovať volebné správanie na základe sociálno-ekonomických podmienok v krajine. 
Dôležité bolo zahrnúť dáta, ktoré zachytávajú stav ekonomiky, kvalitu verejných služieb a životnú úroveň obyvateľstva, pretože práve tieto faktory sú často kľúčovými determinantmi politických preferencií.

Náš dataset teda obsahuje široké spektrum ekonomických a sociálnych ukazovateľov, ktoré nám pomáhajú pochopiť dlhodobé trendy a ich možný dopad na volebné správanie obyvateľstva. 

Medzi najdôležitejšie údaje patrí nezamestnanosť, ktorá odráža stabilitu pracovného trhu a ekonomickú pohodu obyvateľstva. Vysoká nezamestnanosť často vedie k nespokojnosti, ktorá sa môže premietnuť do podpory opozičných strán alebo volania po zmene.
Hrubý domáci produkt (HDP) na obyvateľa poskytuje pohľad na životnú úroveň a ekonomickú produktivitu jednotlivcov, zatiaľ čo celkové HDP ukazuje celkovú ekonomickú silu krajiny. 
Tieto ukazovatele sú dôležité pre pochopenie, ako ekonomická prosperita alebo kríza ovplyvňuje volebné preferencie. 
Index rastu HDP umožňuje sledovať zmenu rastu ekonomiky v čase. Vyššie hodnoty môžu signalizovať optimizmus v spoločnosti, zatiaľ čo pokles môže podporiť nespokojnosť.
Zdravotnícky sektor je reprezentovaný počtom pracovníkov v zdravotníctve. Tento údaj naznačuje dostupnosť a kvalitu zdravotnej starostlivosti, čo má nepriamy vplyv na spokojnosť obyvateľstva a ich dôveru v politické vedenie. 
Výdavky na výskum a vývoj sú kľúčovým ukazovateľom inovačnej kapacity krajiny a schopnosti reagovať na moderné výzvy. Silná podpora výskumu môže naznačovať orientáciu krajiny na budúcnosť a modernizáciu, čo môže ovplyvniť preferencie najmä medzi mladšími voličmi.
Priemerný príjem domácností poskytuje pohľad na materiálnu situáciu obyvateľov. Vyššie príjmy zvyčajne vedú k vyššej spokojnosti a stabilite, čo môže zvýšiť podporu existujúcej vlády, zatiaľ čo nižšie príjmy môžu podporiť volanie po zmene.
Veľkosť investícií do infraštruktúry a verejných služieb naznačuje záväzok vlády zlepšovať životné podmienky občanov. To môže ovplyvniť volebnú podporu, najmä v regiónoch, ktoré čerpajú najviac výhod z týchto projektov.

Tieto dáta sú cenné pri analýze toho, ako sociálne a ekonomické faktory môžu ovplyvniť voľby. Napríklad vyššia nezamestnanosť môže viesť k väčšej podpore strán s populistickým programom, zatiaľ čo ekonomický rast môže zvýšiť dôveru vo vládnuce strany koalície. 
Podobne kvalita zdravotníctva a investície do budúcnosti sú kľúčové pre budovanie dôvery voličov.


\subsection{Election Results}

Nasledne sme si pripravili aj dataset samotných výsledkov volebných období kde sme si strany zoradili a uniformne nazvali.

Dataset o volebných výsledkoch dokumentuje vývoj podpory politických strán v rokoch 2012, 2016, 2020 a 2023. 
Tieto údaje sledujú, ako sa menili politické preferencie obyvateľstva v čase.
Vývoj volebných výsledkov môže byť analyzovaný v kontexte ekonomických a sociálnych ukazovateľov z "General Data" sekcie, čím môžeme identifikovať súvislosti medzi spoločenskými zmenami a politickými preferenciami.

\end{document}